\chapter{Introduction}
\label{ch:intro}

\section{Background}

Secure Shell (SSH) is the standard protocol for remote administration of Unix-like servers,
and it is among the most heavily attacked services on the Internet. A server with port 22
exposed receives password-guessing attempts within minutes of going online; logs such as
the \code{auth.log} reproduced inside our own decoy filesystem are typical. Most of this
traffic comes from automated tools that try common credentials and, on success, run a short
script to download and execute a payload. A smaller share comes from human operators who
explore the system, look for credentials, escalate privileges and install persistence.

Conventional defences (firewalls, rate limiting, intrusion prevention, tools such as
\emph{fail2ban}) are designed to \emph{block} this traffic. Blocking is the right choice for a
production server, but it discards the most informative part of an attack: what the attacker
does after gaining access. Which commands do they run first? Which files do they look for?
Where do they download payloads from, and how do they try to persist?

A \emph{honeypot} answers these questions by inverting the defensive posture. It is a decoy
system with no production value, so any interaction with it is suspicious by definition. It
deliberately admits attackers and records everything they do~\cite{spitzner2002honeypots}.
SSH honeypots such as Kippo and its successor Cowrie~\cite{cowrie} emulate a Linux shell,
log credentials and commands, and capture downloaded files, and they are widely used by
security researchers.

\section{Motivation}

Established SSH honeypots respond \emph{statically}: a given command always produces the same
emulated output, whoever sends it and whatever has happened before. This has two consequences.

\begin{enumerate}
  \item \textbf{Static honeypots cannot make trade-offs.} Sometimes the most informative
    response is the faithful one, because it keeps the attacker engaged. Sometimes a decoy is
    better: a planted credential file can lure an attacker into revealing their next tool. Slowing
    a scripted bot down wastes its time, and disconnecting is reasonable once a session has
    stopped producing information. A fixed policy can make none of these choices.
  \item \textbf{Static honeypots are fingerprintable.} Fixed responses are easy to recognise,
    and an attacker who detects a honeypot stops behaving naturally.
\end{enumerate}

Reinforcement learning (RL) offers a principled way to choose responses: treat each attacker
session as an episode, each response as an action, and the information captured as reward.
Prior work on adaptive honeypots, reviewed in Chapter~\ref{ch:lit}, has shown that this is
feasible. It has rarely asked a quantitative question, though: \emph{how much} better does a
learned policy do than a static one, against which attackers, and compared with what is
achievable at all?

\section{Problem Statement}

\begin{quote}
\emph{Design, implement and evaluate an SSH honeypot whose per-command response is chosen by
a reinforcement-learning policy, with the objective of maximising the quantity and value of
attacker behaviour captured per session, while remaining safe: no attacker input is ever
executed on the host.}
\end{quote}

The project addresses two research questions:

\begin{itemize}
  \item[\textbf{RQ1.}] Does a learned response policy capture more attacker behaviour than a
    static honeypot?
  \item[\textbf{RQ2.}] How does the answer depend on the composition of the attacker population
    and on what the defender values?
\end{itemize}

\section{Objectives}

\begin{enumerate}[label=\textbf{O\arabic*.}]
  \item Build a safe, medium-interaction SSH honeypot that emulates an Ubuntu 22.04 server with
    a per-session in-memory filesystem and over 200 emulated command names.
  \item Formalise honeypot--attacker interaction as a Markov decision process.
  \item Implement an online RL response policy (Q-learning) together with SARSA, bandit,
    static and random baselines that share one interface.
  \item Build an event-driven telemetry pipeline (Redis Streams $\rightarrow$ GeoIP enrichment
    $\rightarrow$ PostgreSQL) with at-least-once delivery.
  \item Build a SOC-style analyst dashboard with session replay and a view of the learned policy.
  \item Deploy all five services reproducibly with Docker Compose, harden the exposed
    component, and measure behaviour under component failure.
  \item Build an attacker simulator and an exact dynamic-programming evaluator, and use them to
    answer RQ1 and RQ2 with benchmarks, confidence intervals and hypothesis tests.
\end{enumerate}

\section{Contributions}

\begin{itemize}
  \item A working five-service system (about 4{,}900 lines of Python and 500 lines of templates, CSS and JavaScript, with 111 automated tests), open
    for inspection at \url{\repourl}.
  \item A formal MDP model of adaptive honeypot responses whose reward measures captured
    intelligence, with a clear account of why the problem is partially observable.
  \item An \emph{exact} evaluation methodology. Each policy's expected return, engagement and
    captured intelligence are computed by solving linear systems on an augmented chain. We also
    compute an oracle upper bound by value iteration and the best policy over the honeypot's own
    observations. Monte-Carlo simulation is used only to validate this machinery.
  \item Empirical findings, including a negative one. Q-learning beats a static honeypot with
    high confidence and recovers 60\% of the available gain, but only with a Robbins--Monro
    step-size schedule. The conventional constant learning rate makes it \emph{worse} than
    static, as stochastic-approximation theory predicts. The benefit scales with the share of
    human attackers.
\end{itemize}

\section{Scope}

\textbf{In scope:} SSH password authentication, interactive and \code{exec} sessions, command
emulation, the response policy, telemetry, the dashboard, containerised deployment and
simulation-based evaluation.
\textbf{Out of scope:} SFTP/SCP file capture, sandboxed analysis of downloaded payloads,
protocols other than SSH, and dashboard authentication (listed as future work).

\section{Organisation of the Report}

Chapter~\ref{ch:lit} reviews honeypots and adaptive honeypots. Chapter~\ref{ch:design} gives
the requirements and system design. Chapter~\ref{ch:formulation} develops the mathematical
formulation: the MDP, the learning algorithm and its convergence conditions, the attacker
model and the exact evaluation method. Chapter~\ref{ch:impl} describes the implementation.
Chapter~\ref{ch:testing} covers testing, security verification and failure experiments.
Chapter~\ref{ch:results} presents the evaluation, and Chapter~\ref{ch:conclusion} concludes.
